\begin{figure}[htbp]
\centering
\begin{mermaid}{
flowchart TB
    classDef node fill:#ffffff,stroke:#555555,stroke-width:1px,color:#000000,font-size:9px

    %% ========== 第2章 ==========
    subgraph CH2[" "]
        direction TB
        T2["<b>第2章  吊钩端雷视融合感知系统设计与标定</b>"]
        subgraph S2_all[" "]
            direction LR
            style S2_all fill:none,stroke:none
            subgraph S2_1["2.1 硬件系统架构与布置"]
                direction TB
                N2_1_1["2.1.1 硬件设备选型"]:::node
                N2_1_2["2.1.2 传感设备安装<br/>布局设计"]:::node
                N2_1_3["2.1.3 吊钩端雷视融合<br/>感知装置设计"]:::node
                N2_1_1 --> N2_1_2 --> N2_1_3
            end
            subgraph S2_2["2.2 通信链路与接口设计"]
                direction TB
                N2_2_1["2.2.1 通信链路<br/>设计与实现"]:::node
                N2_2_2["2.2.2 多传感器时间同步<br/>与延迟补偿"]:::node
                N2_2_1 --> N2_2_2
            end
            subgraph S2_3["2.3 多传感器坐标体系与标定"]
                direction TB
                N2_3_1["2.3.1 坐标系定义与<br/>变换发布约定"]:::node
                N2_3_2["2.3.2 相机内参标定"]:::node
                N2_3_3["2.3.3 双雷达外参标定"]:::node
                N2_3_4["2.3.4 雷达-相机<br/>外参标定"]:::node
                N2_3_1 --> N2_3_2 --> N2_3_3 --> N2_3_4
            end
            S2_1 --- S2_2 --- S2_3
        end
        T2 --- S2_all
    end
    style CH2 fill:#fff3e0,stroke:#ef6c00,stroke-width:2.5px
    style T2  fill:#ef6c00,stroke:#ef6c00,color:#fff

    %% ========== 第3、4章容器 ==========
    subgraph CH34[" "]
        direction LR

        %% ========== 第3章 ==========
        subgraph CH3[" "]
            direction TB
            T3["<b>第3章  基于融合点云的静态场景建图与动态目标感知</b>"]
            subgraph S3_all[" "]
                direction LR
                style S3_all fill:none,stroke:none
                subgraph S3_1["3.1 双雷达点云实时融合与同步方法"]
                    direction TB
                    N3_1_1["3.1.1 双雷达数据接入<br/>与外参应用"]:::node
                    N3_1_2["3.1.2 时间窗口同步与<br/>融合话题发布机制"]:::node
                    N3_1_3["3.1.3 基于IMU的<br/>点云自动调平"]:::node
                    N3_1_1 --> N3_1_2 --> N3_1_3
                end
                subgraph S3_2["3.2 吊钩端位姿估计与静态场景输出"]
                    direction TB
                    N3_2_1["3.2.1 LIO-SAM紧耦合<br/>位姿估计框架"]:::node
                    N3_2_2["3.2.2 双雷达融合定位与<br/>静态场景构建"]:::node
                    N3_2_1 --> N3_2_2
                end
                subgraph S3_3["3.3 动态目标检测与运动状态估计"]
                    direction TB
                    N3_3_1["3.3.1 动态检测输入预处理与<br/>短时历史深度图构建"]:::node
                    N3_3_2["3.3.2 基于遮挡关系的<br/>动态点分类判别"]:::node
                    N3_3_3["3.3.3 欧式聚类候选生成与<br/>短时目标跟踪"]:::node
                    N3_3_1 --> N3_3_2 --> N3_3_3
                end
                S3_1 --- S3_2 --- S3_3
            end
            T3 --- S3_all
        end
        style CH3 fill:#e8f5e9,stroke:#2e7d32,stroke-width:2.5px
        style T3  fill:#2e7d32,stroke:#2e7d32,color:#fff

        %% ========== 第4章 ==========
        subgraph CH4[" "]
            direction TB
            T4["<b>第4章  基于雷视融合的吊物状态估计与三维定位</b>"]
            subgraph S4_all[" "]
                direction LR
                style S4_all fill:none,stroke:none
                subgraph S4_1["4.1 吊装构件视觉识别方法"]
                    direction TB
                    N4_1_1["4.1.1 YOLO26模型<br/>选型与原理"]:::node
                    N4_1_2["4.1.2 吊装构件识别模型<br/>训练与验证"]:::node
                    N4_1_1 --> N4_1_2
                end
                subgraph S4_2["4.2 分割模型量化与边缘部署"]
                    direction TB
                    N4_2_1["4.1.3 分割模型量化与<br/>边缘部署"]:::node
                end
                subgraph S4_3["4.3 雷视跨模态融合与三维定位"]
                    direction TB
                    N4_3_1["4.2.1 基于分割掩码的<br/>点云-图像投影关联"]:::node
                    N4_3_2["4.2.2 吊装构件位置估计与<br/>安全包络建模"]:::node
                    N4_3_1 --> N4_3_2
                end
                S4_1 --- S4_2 --- S4_3
            end
            T4 --- S4_all
        end
        style CH4 fill:#fce4ec,stroke:#c2185b,stroke-width:2.5px
        style T4  fill:#c2185b,stroke:#c2185b,color:#fff

        CH3 --- CH4
    end

    %% ========== 第5章 ==========
    subgraph CH5[" "]
        direction TB
        T5["<b>第5章  吊装作业风险分级预警与避障控制接口设计</b>"]
        subgraph S5_all[" "]
            direction LR
            style S5_all fill:none,stroke:none
            subgraph S5_1["5.1 静态与动态风险分级评估"]
                direction TB
                N5_1_1["5.1.1 基于三维净空的<br/>静态风险评估"]:::node
                N5_1_2["5.1.2 基于相对运动与<br/>碰撞时间的动态风险评估"]:::node
                N5_1_1 --> N5_1_2
            end
            subgraph S5_2["5.2 静动融合的综合风险评估状态机"]
                direction TB
                N5_2_1["5.1.3 静动融合的综合风险<br/>评估状态机设计"]:::node
            end
            subgraph S5_3["5.3 避障控制策略与接口映射"]
                direction TB
                N5_3_1["5.2.1 分层控制<br/>映射框架"]:::node
                N5_3_2["5.2.2 分级动作映射与<br/>避让方向策略"]:::node
                N5_3_3["5.2.3 预警到控制映射<br/>算法及适用边界"]:::node
                N5_3_1 --> N5_3_2 --> N5_3_3
            end
            S5_1 --- S5_2 --- S5_3
        end
        T5 --- S5_all
    end
    style CH5 fill:#f3e5f5,stroke:#7b1fa2,stroke-width:2.5px
    style T5  fill:#7b1fa2,stroke:#7b1fa2,color:#fff

    %% ========== 跨章连接 ==========
    CH2 --> CH34
    CH34 --> CH5
}
\end{mermaid}
\caption{技术路线图}
\label{fig:tech_route}
\end{figure}
